\documentclass[11pt]{article}
\usepackage[T1]{fontenc}
\usepackage{lmodern,amsmath,amssymb,amsthm,mathtools}
\usepackage[margin=1in]{geometry}
\usepackage{microtype,booktabs,longtable,array,enumitem,xcolor}
\usepackage{fancyhdr}
\usepackage[hidelinks]{hyperref}
\usepackage{xurl}
\hypersetup{pdftitle={Prime-supported Pell growth and binomial sequence spacing},pdfauthor={AI-assisted research prepared for Vladimir Reshetnikov}}
\setlength{\headheight}{14pt}
\pagestyle{fancy}
\fancyhf{}
\fancyhead[L]{\small Prime-supported spacing for Diophantine consistency}
\fancyhead[R]{\small 25 October 2026}
\fancyfoot[C]{\thepage}
\setlength{\parskip}{0.35em}
\setlength{\parindent}{0pt}
\setlist{nosep,leftmargin=*}
\newtheorem{theorem}{Theorem}[section]
\newtheorem{lemma}[theorem]{Lemma}
\newtheorem{proposition}[theorem]{Proposition}
\newtheorem{corollary}[theorem]{Corollary}
\theoremstyle{definition}
\newtheorem{definition}[theorem]{Definition}
\theoremstyle{remark}
\newtheorem{remark}[theorem]{Remark}
\newcommand{\Z}{\mathbb Z}
\newcommand{\N}{\mathbb N}
\newcommand{\ZFC}{\mathsf{ZFC}}
\newcommand{\Con}{\operatorname{Con}}
\newcommand{\lcm}{\operatorname{lcm}}
\newcommand{\vval}{\nu}
\newcommand{\Ccal}{\mathcal C_{222}}
\newcommand{\NumVariables}{209}
\newcommand{\NumRows}{150}
\newcommand{\NumMonomials}{5283}
\newcommand{\PolyHeight}{274899927541}
\newcommand{\RowHeight}{524309}
\newcommand{\GrowthPrime}{524309}
\newcommand{\NumResidualTerms}{999}

\title{Prime-supported Pell growth\\and binomial sequence spacing\\[0.4em]\large An explicit quartic in 209 unrestricted integers\\whose unsolvability is equivalent to $\Con(\ZFC)$}
\author{Research prepared for Vladimir Reshetnikov and ProveIt\\\small AI-assisted construction, proof, and independent arithmetic review}
\date{25 October 2026}
\begin{document}
\maketitle
\thispagestyle{empty}
\begin{abstract}
We give a fully specified polynomial $P\in\Z[x_0,\ldots,x_{208}]$ of
degree four and an ordinary proof, using the explicitly identified inherited
calibration of the numerical calculus $\Ccal$, that ZFC proves
$\neg\exists\boldsymbol x\in\Z^{209}\,P(\boldsymbol x)=0$
equivalent to $\Con(\ZFC)$. The polynomial is a sum of 150 squares of
quadratic residuals; all residuals are printed in this report. Its collected
expansion has 5,283 nonzero monomials and coefficient height
274,899,927,541.

The new construction couples two previously separate witnesses. A Pell
ordinate supplies both a large growth scale and divisibility by every small
prime. An already present binomial selector then supplies the capacity needed
for the spacing of a G\"odel sequence code. This removes the independent
spacing multiplier. The key completeness argument chooses the proof horizon
from the cofinal family $t=32c^{2k}$, with the fixed prime $c=524309$;
finite-field bounds on Pell annihilating indices make the necessary
congruences compatible.

Relative to the frozen incoming 210-variable quartic, the construction saves
one unrestricted integer variable and reduces coefficient height by almost
a factor of four, while adding one collected monomial. It does not establish
a globally simplest equation or worldwide priority. The unbounded proof is
distinct from the supplied exact coefficient checks. No new end-to-end
Lean/Rocq verification or proof of consistency is claimed.
\end{abstract}
\clearpage
\tableofcontents
\clearpage

\section{Result, comparison, and meaning of simplicity}

The primary result is the concrete equation
\begin{equation}\label{eq:main}
 \boxed{\quad P(x_0,\ldots,x_{208})
       :=\sum_{j=0}^{149}R_j(x_0,\ldots,x_{208})^2=0,
       \qquad \boldsymbol x\in\Z^{209}.\quad}
\end{equation}
Appendix~\ref{app:residuals} specifies every $R_j$ as a literal polynomial
with integer coefficients. These 209 coordinates are all unrestricted
integers and all occur in the collected polynomial. Auxiliary letters used
in the mathematical explanation are abbreviations or proof parameters;
they do not enlarge the variable list in~\eqref{eq:main}.

\begin{theorem}[Main result]\label{thm:main}
For the literal polynomial in Appendix~\ref{app:residuals}, using the
retained ordinary arithmetic and $\Ccal$/ZF/ZFC calibration described in
Section~\ref{sec:inherited}, there is a proof formalizable in ZFC of
\[
 \bigl(\exists\boldsymbol x\in\Z^{209}\ P(\boldsymbol x)=0\bigr)
 \quad\longleftrightarrow\quad
 \exists p\,\operatorname{Proof}_{\ZFC}(p,\bot).
\]
Consequently,
\[
 \ZFC\ \vdash\
 \bigl(\neg\exists\boldsymbol x\in\Z^{209}\ P(\boldsymbol x)=0\bigr)
 \longleftrightarrow\Con(\ZFC).
\]
\end{theorem}

Here $\operatorname{Proof}_{\ZFC}$ refers to the fixed effective presentation
used by the inherited calibration. The assertion is internal: the proof
codes and integers are those quantified over in ZFC. No restriction to
externally standard codes, transitive models, or externally well-founded
models is added. The equivalence does not assume consistency or
inconsistency, and does not decide which side holds.

The immediate comparison is the incoming package
\texttt{ConZFC\_Pell\_Growth\_210\_2026-10-25.zip} at the pinned repository
commit~\cite{parent,repo}. Its complete bytes are retained in this release.
Both columns below were independently collected from their literal
residual lists.

\begin{center}
\begin{tabular}{lrr}
\toprule
Measure & Frozen parent & This report\\
\midrule
Unrestricted integer variables & 210 & \textbf{209}\\
Quadratic residual squares & 150 & 150\\
Exact total degree & 4 & 4\\
Collected monomials & 5,282 & 5,283\\
Collected coefficient height & 1,099,511,627,776 & 274,899,927,541\\
Largest residual coefficient & 1,048,576 & 524,309\\
Residual monomials before squaring & 999 & 999\\
\bottomrule
\end{tabular}
\end{center}

The new height is about $25.0020\%$ of the old one. The single extra
collected monomial is an explicit cost. We report the expanded coefficient
height, not just the largest scalar in a circuit.

There is no single established ordering called ``simplest'' that combines
variable count, degree, coefficient height, support, circuit size, domain,
and proof complexity. The inspected repository already contains equations
with many fewer variables at other degrees and coefficient costs, including
the incoming fourteen-positive-variable, 65-operation carrier equation
\cite{carrier}. This work improves the direct quartic family. It uses the
numerical proof calculus, without a K215 simulation or an instantiation of
Jones's universal polynomial. It does not show that 209 is a minimum, or
that degree four is necessary for this particular consistency statement.

\section{The new idea in arithmetic terms}\label{sec:idea}

The direct compiler needs two kinds of size. First, $F$ must be large enough
to separate products of formula weights $1+Fj$. Second, a spacing $y$ must
be large enough to store the intermediate pending-multiset products in a
G\"odel $\beta$ code. For completeness the moduli
\[
 1+y,\ 1+2y,\ \ldots,\ 1+(t+1)y
\]
must be pairwise coprime. A sufficient condition is that every prime at most
$t$ divide $y$.

The parent uses a freely chosen odd multiplier:
\[
 y=F(2v+1).
\]
It chooses $v$ sufficiently large and with appropriate divisibility. We
replace this by
\begin{equation}\label{eq:spacing}
 \boxed{\ y=FB,\qquad B=\binom{F+t}{t}.\ }
\end{equation}
The coordinate $B$ and its exact polynomial authentication already exist.
The new equation for $y$ is quadratic and the coordinate $v$ disappears.
The binomial notation occurs in the proof, not as a primitive operation in
the final polynomial.

This substitution has two independent obligations. The value $FB$ must
exceed every state of the selected finite proof. Also $FB$ must have enough
small prime divisors for the beta moduli. The binomial supplies the first
property through its degree in $F$. We arrange the second directly in $F$.

Put
\begin{equation}\label{eq:constants}
 d=2^{19}=524288,\qquad c=524309.
\end{equation}
The constant $c$ is prime: exact trial division by every prime not exceeding
$\lfloor\sqrt c\rfloor=724$ verifies this finite fact. Impose
\begin{align}
 t&=32(2z+1)^2,\label{eq:horizon}\\
 Z&=ctF,\label{eq:g1}\\
 X^2-Z^2-2ZF&=1,\label{eq:g2}\\
 F+t&=ctu.\label{eq:g3}
\end{align}
All four residuals have degree at most two. The middle pair gives
\begin{equation}\label{eq:pell}
 X^2-\bigl((ct+1)^2-1\bigr)F^2=1.
\end{equation}
The congruence $F\equiv-t\pmod{ct}$ will force a sufficiently large
positive Pell index, once the existing signed arithmetic front has
established $F>0$.

For completeness we may select horizons from the particular cofinal family
\begin{equation}\label{eq:cofinal}
 t=32c^{2k},\qquad z=(c^k-1)/2,\qquad k\ge1.
\end{equation}
An arbitrary zero of the final polynomial need not have this special form.
It only satisfies~\eqref{eq:horizon}. The stronger choice is made when
constructing a zero from a finite refutation. This distinction is essential.

The role of the prime $c$ is structural. The modulus $ct$ contains one more
factor of $c$ than $t$ does, but every small-prime Pell annihilating index
has fewer than $2k+1$ factors of $c$. Therefore its required divisibility
can coexist with the residue $n\equiv-t\pmod{ct}$. This is the new
compatibility argument that lets the growth witness absorb the prime
support formerly delegated to $v$.

\section{What is inherited and what is proved here}\label{sec:inherited}

The retained parent package contains the literal polynomial source, exact
integer projections, arithmetic proofs, the relevant calibration sources,
and its own earlier source archive. It is cited as a proved predecessor,
with the following boundaries made explicit.

\subsection{Arithmetic interfaces retained from the predecessor}

The unchanged rows provide the following facts and constructions
\cite{parent,anchor,linearcrt}.
\begin{enumerate}
\item A shared signed Pell frame authenticates the retained power requests.
The principal pairs are oriented before the signs of the arithmetic front
are inferred. At any authenticated base of absolute value at least two,
an integral power output forces its nominal exponent to be nonnegative and
equals the intended power.
\item The actual CRT index representative $C$ is even and nonzero. Its
Pell pair is the reused anchor. The cap bounds all other retained nominal
indices strictly below $C$ after nonnegativity is recovered.
\item The arithmetic front retains
\begin{align}
 B\eta&=C+F+1,& D&=2(C+F+1),& q&=D^x,\label{eq:front1}\\
 x&=F+1+(y+1)r_0,&
 x&=1+\bigl(1+(t+1)y\bigr)r_t.\label{eq:endpoints}
\end{align}
The powers with base $q$ and exponent $y+1$, and with base $q+1$ and
exponents $C$ and $F+t$, are among the retained requests.
\item Exact digit extraction authenticates $B=\binom{F+t}{t}$ and the
range-binomial predicates. The shared remainder bound and the affine radix
$D$ suffice because all relevant upper indices are at most $C$ and
$q=D^x>(2C)^x$.
\item The factorial-free binomial CRT construction encodes the local
calculus constraints. Its completeness requires the established growth
bound on $F$, an odd selector $B$, a suitable positive beta spacing and
cofinal choices of the CRT representatives.
\end{enumerate}

We do not delete or weaken an anchor, orientation, cap, range, or remainder
row. The new $X,Z,u$ are private growth coordinates; they are not new
requested atoms of the shared power frame. The old actual front atoms
$F,B,t,y+1$ are already present in that frame.

\subsection{Logical calibration}

The direct arithmetic source uses an inherited square-shell recoding of
the numerical calculus $\Ccal$. Its syntax translation fixes truth code 0
and falsity code 1 and transports every guarded rule instance. The displayed
square-shell constructors are not literally the original Lean constructors.
Compose this recoding with the retained original-calculus interfaces
\nolinkurl{KEndpoint.c222_iff_zf_refutation} and
\nolinkurl{ConstructibleChoice.zf_refutation_iff_zfc_refutation}
\cite{anchor}. Their ordinary proofs supply the connection with the fixed
ZFC presentation. The first interface is a statement about the numerical
calculus itself; the occurrence of ``KEndpoint'' in a file or namespace
does not insert a Turing-machine simulation into this construction.

The source includes a depth certificate for the square-shell constructor
trees, and a Horn-flow proof extracting a finite refutation from the
authenticated local arithmetic. These are inherited mathematical inputs.
This report proves the new growth, prime-support, capacity, sign, and
completeness arguments in full. It does not reprove the entire logical
calibration or package a complete fresh Lean dependency closure.

Thus the result is an ordinary mathematical proof using named predecessor
lemmas, together with exact finite verification of the delivered polynomial.
It is not a newly checked object-language ZFC derivation. The successful
Python checks are not a substitute for the unbounded arithmetic or logical
theorems.

\section{Elementary Pell facts}\label{sec:pellfacts}

For an integer $a\ge2$, define
\[
 X_0(a)=1,\quad Y_0(a)=0,
\]
and
\begin{align*}
 X_{n+1}(a)&=aX_n(a)+(a^2-1)Y_n(a),\\
 Y_{n+1}(a)&=X_n(a)+aY_n(a).
\end{align*}
Equivalently,
\[
 X_n(a)+Y_n(a)\sqrt{a^2-1}
   =\bigl(a+\sqrt{a^2-1}\bigr)^n.
\]
The recurrence, rather than real exponentiation, suffices for all subsequent
formalization claims.

\begin{lemma}[Classification and elementary bounds]\label{lem:pell}
For $a\ge2$, every integral solution of
$X^2-(a^2-1)Y^2=1$ with $Y\ne0$ has
$|X|=X_n(a)$ and $|Y|=Y_n(a)$ for some positive integer $n$.
Moreover,
\begin{align}
 Y_n(a)&\equiv n\pmod{a-1},\label{eq:baseone}\\
 Y_n(a)&\ge(2a-1)^{n-1}\quad(n\ge1).\label{eq:pellgrowth}
\end{align}
If $a$ is odd, then
\begin{equation}\label{eq:v2}
 \vval_2(Y_n(a))=\vval_2(n)\quad(n\ge1).
\end{equation}
\end{lemma}
\begin{proof}
For classification take $x=|X|$, $y=|Y|>0$. If $y=1$, then $x=a$.
In general $(a-1)y<x\le ay$, and multiplying $x+y\sqrt{a^2-1}$ by
$a-\sqrt{a^2-1}$ gives another nonnegative integral norm-one pair with
ordinate $ay-x$ strictly smaller than $y$. Descent reaches $(1,0)$.
Reversing the descent proves classification and uniqueness of the positive
index.

The recurrence defines integer polynomials in $a$. At $a=1$ it gives
$X_n(1)=1$ and $Y_n(1)=n$, proving~\eqref{eq:baseone}. The second-order
recurrence
$Y_{n+1}=2aY_n-Y_{n-1}$ and strict increase imply
$Y_{n+1}\ge(2a-1)Y_n$; use $Y_1=1$ for~\eqref{eq:pellgrowth}.

When $a$ is odd, the recurrence modulo two makes every $X_n$ odd and
$Y_n\equiv n\pmod2$. The doubling identity
$Y_{2m}=2X_mY_m$ then proves~\eqref{eq:v2} by writing $n=2^r s$ with
$s$ odd.
\end{proof}

\begin{lemma}[A bounded annihilating index]\label{lem:annihilator}
Let $a\ge2$, $\delta=a^2-1$, and let $p$ be an odd prime. Set
\[
 \rho_p=
 \begin{cases}
  p,&p\mid\delta,\\
  p-\left(\dfrac{\delta}{p}\right),&p\nmid\delta,
 \end{cases}
\]
where the second line uses the Legendre symbol. Then
$1\le\rho_p\le p+1$ and $p\mid Y_{\rho_p}(a)$. Any positive multiple
of $\rho_p$ has the same vanishing property. If $a$ is odd, we may take
$\rho_2=2$.
\end{lemma}
\begin{proof}
If $p\mid\delta$, the recurrence modulo $p$ gives
$Y_n(a)\equiv n a^{n-1}$ and $a\equiv\pm1$, so $n=p$ works.
Otherwise let $s^2=\delta$ in $\mathbb F_p$ when $\delta$ is a square,
or in $\mathbb F_{p^2}$ when it is a nonsquare. Put $w=a+s$; its
inverse is $a-s$. In the split case $w^{p-1}=1$. In the nonsplit case
Frobenius sends $s$ to $-s$, so $w^p=w^{-1}$ and $w^{p+1}=1$.
In either case
\[
 Y_n(a)=\frac{w^n-w^{-n}}{2s}
\]
vanishes at the stated index. The denominator is nonzero. The addition
identities, or powers of $w$, give vanishing at all multiples.
For $p=2$ and odd $a$, use $Y_2(a)=2a$ and the addition identities.
\end{proof}

The lemma does not assert that $\rho_p$ is the smallest possible index.
Only its explicit upper bound is needed; no theorem about minimal ranks of
apparition is required.

\section{Growth and the signed arithmetic front}\label{sec:soundness}

\subsection{The inherited label bound}

The square-shell constructors are
\begin{align*}
 I(u,v)&=2\bigl((u+v)^2+u\bigr),& M(u,v)&=I(u,v)+1,\\
 A(u,v)&=2\bigl((u+v)^2+3u+4v+3\bigr),& E(u,v)&=A(u,v)+1.
\end{align*}
For $0\le u,v\le M_0$, each constructor value plus one is at most
\[
 E(M_0,M_0)+1=8M_0^2+14M_0+8\le8(M_0+1)^2.
\]
With each local parameter below $t$, a depth-$r$ tree therefore has value
plus one at most $8^{2^r-1}t^{2^r}$. The retained source certificate has
maximum constructor depth 19. Hence its local labels satisfy
\begin{equation}\label{eq:labelbound}
 0\le j<L,\qquad L=8^{d-1}t^d,\qquad d=2^{19}.
\end{equation}
This is a bound on decoded local values. It is not an upper bound on their
large CRT representatives.

\subsection{A sign-free bound comes first}

By~\eqref{eq:horizon}, $t\ge32$ and $t$ is even. Equation~\eqref{eq:g3}
says $F=t(cu-1)$, so $F\ne0$ because $c\nmid1$, and $F$ is even.
Apply Lemma~\ref{lem:pell} at $a=ct+1$. There is $n\ge1$ with
$|F|=Y_n(a)$. The congruence~\eqref{eq:baseone} and
$F\equiv-t\pmod{ct}$ imply
\[
 n\equiv t\text{ or }-t\pmod{ct}.
\]
Both least positive residues are at least $t$. Thus
\begin{equation}\label{eq:weakgrowth}
 |F|\ge(2ct+1)^{t-1}>t.
\end{equation}
No assertion about the sign of $F$ has been used.

The shared frame establishes its canonical-index conclusions independently
of these new private coordinates. Since $C$ and $F$ are even,
$C+F+1$ is odd. The selector relation in~\eqref{eq:front1} makes $B$
and $\eta$ odd and nonzero. Therefore the new relation $y=FB$ gives
\begin{equation}\label{eq:weakspacing}
 |y|\ge|F|>t.
\end{equation}

\subsection{Recovering signs without a circular assumption}

Here is the complete order needed from the retained power and endpoint
interfaces.
\begin{enumerate}
\item $D=2(C+F+1)$ is even and nonzero, so $|D|\ge2$. The request
$q=D^x$ implies $x\ge0$.
\item If $x=0$, the second endpoint in~\eqref{eq:endpoints} makes
$1+(t+1)y$ divide $-1$. But its absolute value exceeds one
by~\eqref{eq:weakspacing}. Hence $x\ge1$ and $|q|\ge2$.
\item The request at base $q$ and exponent $y+1$ forces $y+1\ge0$.
A negative $y$ would have $y\le-|F|<-t$, which is impossible.
Thus $y\ge|F|>0$.
\item If $x=1$, the first endpoint gives $-F=(y+1)r_0$, contrary to
$0<|F|\le y<y+1$. Hence $x\ge2$.
The second endpoint now yields $r_t\ge1$ and
\begin{equation}\label{eq:xbound}
 x\ge2+(t+1)y.
\end{equation}
\item Consequently $|q|\ge4$ and $|q+1|\ge3$. The retained requests
at base $q+1$ force their nominal indices nonnegative. In particular
$F+t\ge0$. Together with~\eqref{eq:weakgrowth} this gives $F>0$.
The same argument gives $C\ge0$; the anchor already excluded $C=0$.
Thus $C>0$, $D>0$, and $q>0$.
\item Since $y=FB$ and $y,F>0$, also $B>0$. The first endpoint gives
$r_0\ge0$ and $x\ge F+1$.
\end{enumerate}

Because $B$ is odd, the old spacing multiplier could now be reconstructed as
$v=(B-1)/2\in\Z$ if needed to invoke a retained beta lemma. It is not a
quantified coordinate of the new polynomial. All sign and remainder
arguments used by the old arithmetic front therefore retain their premises.

\subsection{Recovering the full separation budget}

Now $F>0$, so the ambiguous congruence resolves to
\[
 n\equiv-t\pmod{ct},\qquad n\ge(c-1)t>dt.
\]
Since $2c\ge8$, Lemma~\ref{lem:pell} gives
\begin{align*}
 F&\ge(2ct+1)^{(c-1)t-1}
   \ge(8t)^{dt-1},\\
 \frac{(8t)^{dt-1}}{L^t}
   &=\frac{8^{t-1}}t\ge1.
\end{align*}
Therefore
\begin{equation}\label{eq:fullgrowth}
 F\ge L^t.
\end{equation}
This is precisely the inherited product-separation hypothesis. It is
derived after the sign bootstrap, not used to justify that bootstrap.

The unchanged affine-radix and digit arguments now authenticate the binomial
selector and local range predicates. For example, $D>2C$ and the cap imply
that every required binomial coefficient is bounded by $C^x<q/2^x$.
The selector relation also gives $|B|\le D/2<q/2$ before its digit is
identified. The retained common remainder bounds select the exact digit.
The inherited separation and Horn-flow arguments then recover a finite
$\Ccal$ refutation. This proves soundness of the modified source.

\newpage
\section{The prime-support completeness theorem}\label{sec:prime}

\begin{theorem}[Cofinal Pell ordinates with small-prime support]
\label{thm:primesupport}
Let $c=524309$ and $t=32c^{2k}$ with $k\ge1$. Put $a=ct+1$.
For any integer $\ell\ge1$, there are arbitrarily large positive integers
$n$ satisfying
\[
 n\equiv-t\pmod{ct},\qquad 2^\ell\mid Y_n(a),\qquad
 p\mid Y_n(a)\quad\text{for every prime }p\le t.
\]
Consequently the new growth equations have cofinally large positive
solutions $F$ with all those divisibility properties.
\end{theorem}
\begin{proof}
For each prime $p\le t$ choose the explicit $\rho_p$ of
Lemma~\ref{lem:annihilator}, including $\rho_2=2$. Set
\[
 M=\lcm\bigl(2^\ell,\{\rho_p:p\le t\}\bigr).
\]
Since $c>33$,
\[
 \rho_p\le p+1\le t+1=32c^{2k}+1<c^{2k+1}.
\]
Thus $\vval_c(\rho_p)\le2k$, and $\vval_c(M)\le2k$.
The only prime factors of $ct=2^5c^{2k+1}$ are 2 and $c$. Its
intersection with $M$ has 2-adic valuation at most five and $c$-adic
valuation at most $2k$. Therefore
\begin{equation}\label{eq:crtcompat}
 \gcd(M,ct)\mid 2^5c^{2k}=t.
\end{equation}
The two congruences
\[
 n\equiv0\pmod M,\qquad n\equiv-t\pmod{ct}
\]
are compatible by the generalized Chinese remainder theorem. Their
solutions form an arithmetic progression with positive period, so positive
solutions are cofinal.

For such $n$, every $\rho_p$ divides $n$; Lemma~\ref{lem:annihilator}
gives $p\mid Y_n(a)$. Also $a$ is odd and $2^\ell\mid n$, so
Lemma~\ref{lem:pell} gives $2^\ell\mid Y_n(a)$. Put
\[
 F=Y_n(a),\quad X=X_n(a),\quad Z=ctF,\quad
 u=\frac{F+t}{ct}.
\]
The last quotient is integral by~\eqref{eq:baseone}; all three growth
rows hold. Strict increase of $Y_n(a)$ makes the resulting positive
ordinates cofinal as well.
\end{proof}

\begin{corollary}[Odd selector and coprime beta moduli]\label{cor:beta}
Choose $\ell$ with $2^\ell>t$, take a positive $F$ from
Theorem~\ref{thm:primesupport}, and define
\[
 B=\binom{F+t}{t},\qquad y=FB.
\]
Then $B$ is odd and the integers $1+iy$, $1\le i\le t+1$, are
pairwise coprime.
\end{corollary}
\begin{proof}
For $1\le i\le t<2^\ell$, divisibility $2^\ell\mid F$ implies
$\vval_2(F+i)=\vval_2(i)$. Subtracting valuations in
\[
 B=\frac{\prod_{i=1}^{t}(F+i)}{t!}
\]
gives $\vval_2(B)=0$.

Every prime at most $t$ divides $F$, and hence $y$. If a prime $p$
divided both $1+iy$ and $1+jy$ for $i<j$, it would be coprime to $y$
and divide $j-i$. Thus $p\le j-i\le t$, contradicting $p\mid y$.
\end{proof}

The proof requires only prime support in $y$, not $t!\mid y$. It also
does not require every solution of the new polynomial to have the
small-prime divisibility properties. Those are choices in the completeness
construction; soundness already follows from the retained exact arithmetic.

\section{Why the binomial selector has enough capacity}\label{sec:capacity}

The number of available proof rows can be chosen strictly larger than the
number of genuine backward-expansion steps. Starting from one pending
formula, each step replaces one formula by at most two premises. A proof
with $h$ such steps therefore has at most $h+1$ pending formulas at any
time. The retained truth padding leaves the pending product unchanged
\cite{parent}. Choose the horizon $t$ with $t>h+1$ and above all finite
local parameter requirements. The family~\eqref{eq:cofinal} permits this.

\begin{lemma}[Binomial capacity]\label{lem:capacity}
Let $t\ge2$, $L\ge t$, and $F\ge L^t$ be integers. Suppose a state is
a product of at most $t-1$ weights $1+Fj$, where $0\le j<L$. Then
\[
 S<F\binom{F+t}{t}.
\]
\end{lemma}
\begin{proof}
Since the labels are integers,
$1+Fj\le1+F(L-1)\le FL$. Thus
$S\le(FL)^{t-1}$, including the empty product. On the other hand,
\[
 F\binom{F+t}{t}\ge\frac{F^{t+1}}{t!}.
\]
The ratio of the latter lower bound to the former upper bound is
\[
 \frac{F^2}{t!L^{t-1}}
 \ge\frac{L^{2t}}{L^t L^{t-1}}=L>1,
\]
where $t!\le t^t\le L^t$. This proves the strict inequality.
\end{proof}

Apply the lemma to~\eqref{eq:labelbound} and~\eqref{eq:fullgrowth}.
Every actual pending state is below $y=FB$, so it is a valid remainder
for each relevant beta modulus. No additional growth request, inequality
coordinate, or free spacing multiplier is needed.

There is also a useful alternative proof of the existence needed here.
For any fixed proof and horizon, its pending states are polynomials in $F$
of degree at most $t-1$, while $F\binom{F+t}{t}$ is a polynomial of
degree $t+1$ with positive leading coefficient. The cofinal Pell
progression eventually dominates them. The explicit estimate above is
stronger: the already authenticated lower bound on $F$ suffices.

\section{Completeness and the internal consistency equivalence}
\label{sec:complete}

Start with a finite $\Ccal$ refutation. Unfold it to a finite backward
proof schedule if necessary, and let $h$ be its number of genuine expansion
steps. This may enlarge the proof, but the construction has no imposed
resource bound.

Choose $k$ sufficiently large that $t=32c^{2k}$ is greater than $h+1$
and all required finite local parameters. Then set $z=(c^k-1)/2$.
Choose $\ell$ with $2^\ell>t$ and use
Theorem~\ref{thm:primesupport} to construct $n,F,X,Z,u$.
Define the already intended selector and spacing by
\[
 B=\binom{F+t}{t},\qquad y=FB.
\]
Corollary~\ref{cor:beta} makes the beta moduli pairwise coprime and $B$
odd. Lemma~\ref{lem:capacity} bounds every required pending state. Truth
steps pad the proof to the selected horizon without enlarging those states.

Ordinary finite CRT now supplies a positive beta code $x$ with the
prescribed states and endpoints, and permits $x$ to be arbitrarily large
in its residue class. Set $N=B^x$. The unchanged factorial-free
binomial CRT construction provides the residue classes for the ordinary
representatives and $C$ at precision $N^2$. The growth hypothesis is
stronger than its lower-bound requirement.

Because $N$ is odd, the prescribed class for $C$ has arbitrarily large
even positive representatives. Choose the other representatives first,
then choose $C$ even and larger than all other retained nominal indices,
including $F+t$, $x$, and $y+1$. The retained selector congruence gives
$B\mid C+F+1$. Set
\[
 \eta=(C+F+1)/B,\qquad D=2(C+F+1),\qquad q=D^x.
\]
The same binomial-digit bounds and $N\mid q$ properties as in the
predecessor follow. Choose the common frame center and the canonical
principal Pell pairs last. The inherited cap, orientation, power,
remainder, and local arithmetic completeness constructions supply all
remaining coordinates.

The order is explicit and acyclic: the finite proof fixes $t$; then the
growth construction fixes $F,B,y$; then beta CRT fixes $x$ and $N$; then
the local CRT representatives and $C$ are chosen; finally $D,q$ and the
shared frame are completed. No modulus is chosen using a coordinate that
was required to have been fixed earlier.

Soundness was established in Section~\ref{sec:soundness}; completeness is
the construction just given. Compose them with the retained
$\Ccal\leftrightarrow\text{ZF refutation}\leftrightarrow
\text{ZFC refutation}$ calibration to obtain Theorem~\ref{thm:main}.
Every new step is finite arithmetic, induction, CRT, or elementary finite
field algebra, and is formalizable in ZFC. The cofinal choices are proved
existence claims inside that same theory.

\section{From the semantic change to the literal polynomial}
\label{sec:literal}

The retained parent equation has 210 coordinates and 150 residuals.
Its exact SHA-256 is
\begin{quote}\small\ttfamily
\nolinkurl{20a2108353a3756ed7d1ddfdc71695131dcf0b6d94039988d9565ac48e6e1f0a}.
\end{quote}
The new equation is obtained by the following finite specification.
All indices in this list are \emph{parent} indices until the final
renumbering.
\begin{enumerate}
\item Undo the parent's optimization $R_3\gets R_3+R_{146}$ by
replacing row 3 with $R_3-R_{146}$. This is an invertible integer row
operation and restores the literal selector relation before beta is changed.
\item Replace row 125 by
$x_0-32-128x_{178}-128x_{178}^2$.
\item Replace row 146 by $x_4-1-x_2x_3$.
\item Replace rows 147--149 by
\begin{align*}
 x_{208}-c x_0x_2,&\\
 x_{207}^2-x_{208}^2-2x_2x_{208}-1,&\\
 x_2+x_0-cx_0x_{209}.&
\end{align*}
\item Delete the now-unused coordinate 206. Keep coordinates 0--205
unchanged and send old coordinates 207, 208, 209 to new coordinates
206, 207, 208.
\end{enumerate}
No other residual is changed. Exactly the following semantic coordinates
are relevant to the modification.

\begin{center}
\begin{tabular}{lll}
\toprule
Quantity & Parent coordinate & Final coordinate\\
\midrule
$t$ & $x_0$ & $x_0$\\
$F$ & $x_2$ & $x_2$\\
$B$ & $x_3$ & $x_3$\\
$y+1$ & $x_4$ & $x_4$\\
$C$ & $2x_{16}$ & $2x_{16}$\\
$z$ & $x_{178}$ & $x_{178}$\\
Old free multiplier & $x_{206}$ & deleted\\
$X$ & $x_{207}$ & $x_{206}$\\
$Z$ & $x_{208}$ & $x_{207}$\\
$u$ & $x_{209}$ & $x_{208}$\\
\bottomrule
\end{tabular}
\end{center}

The new equation JSON has SHA-256
\begin{quote}\small\ttfamily
\nolinkurl{b0413e587c595c6813f8eb67ef1581a507cee6155a38b8d091ca064f3e29b9f6}.
\end{quote}
The separate collected expansion fixes all 5,283 coefficients. In its
term format, $[a,[i,j,k]]$ means $a x_ix_jx_k$; repeated indices mean
powers and an empty list means a constant. The equation file fixes the
residual squares instead, and is substantially easier to audit.

The semantic replacements are not asserted to be invertible maps on all
parent zeros. In particular an arbitrary old refutation witness need not
already have prime-supported $F$. The completeness proof constructs an
appropriate new witness from the underlying finite refutation.
Only the initial row operation and final unused-coordinate deletion are
pure algebraic operations of the elementary kind stated above.

\section{Verification, reproducibility, and scope}\label{sec:verify}

The release provides separate implementations for production and finite
verification. The verifier does not import the producer's polynomial code
or the predecessor's polynomial helpers. It authenticates the parent
bytes, reconstructs the six changed rows, checks the renumbering and
variable occurrence, independently collects every square, and compares
the full expansion and metrics. It also verifies the fixed prime by trial
division. Its checks remain active with Python optimization enabled.

Normal and \texttt{-O} runs pass. Seven deliberate residual corruptions
and one collected-coefficient corruption are rejected. A separate parser
checks every one of the 150 printed residuals and all 999 printed
coefficients against the equation JSON, with a deliberate print-corruption
control. The parent package's complete release verifier was freshly run;
its 352-file manifest, inherited checks, source/projection chain, printed
equations, and fresh reconstruction passed. The receipt records the
actual scope rather than relabeling an old receipt as a new run.

The arithmetic regression program checks 1,248 local Pell annihilating
cases, including 126 ramified cases. A smaller-prime completeness fixture
uses $c=37$, $t=43,808$, 4,564 prime constraints, and a 10,420-bit index.
It checks the congruences by modular powering rather than attempting to
print the enormous Pell ordinate. These are finite debugging checks;
the proof for all horizons is Theorem~\ref{thm:primesupport}.

Two independently assigned AI review tasks checked the new mathematical
arguments. Their notes are retained. This is not external human peer review.
No numerical search for a root can validate the consistency equivalence,
and none is presented as doing so.

From the extracted package root, Python 3.10 or later suffices for the
finite checks:
\begin{verbatim}
python3 -B code/verify.py
python3 -B -O code/verify.py
python3 -B code/check_printed.py
python3 -B code/check_arithmetic.py
\end{verbatim}
The literal equation can be regenerated with
\begin{verbatim}
python3 -B code/build.py
\end{verbatim}
The evaluator accepts a JSON array containing exactly 209 integers:
\begin{verbatim}
python3 -B code/evaluate.py tuple.json --residuals
\end{verbatim}
To rebuild the standalone article and PDF with a standard TeX Live
installation, run
\begin{verbatim}
python3 -B code/build_report.py
\end{verbatim}
PDF metadata or TeX versions can change PDF bytes. The finite equation
identity checks are independent of such changes. \texttt{SHA256SUMS}
authenticates the delivered files; it is not a theorem about the arithmetic.

\section{What this adds, and the remaining research problem}

The contribution is the proved reuse of a selector for sequence spacing,
made possible by a cofinal Pell construction with controlled prime support.
It couples two arithmetic tasks that were independent in the parent.
The prime-supported ordinate supplies the coprimality requirement; the
binomial selector supplies capacity; the existing front supplies sign
recovery and the exact selector interpretation. None of these roles is
left to an unspecified universal polynomial or an uncompiled variable
exponent.

The quantitative gain is modest and exact: one fewer unrestricted
integer unknown, unchanged degree and residual count, nearly four times
smaller coefficient height, and one additional collected monomial. The
argument supplies a reusable arithmetic lemma, which may have value beyond
this one-coordinate saving. It does not establish that its proof complexity
is smaller than the parent's.

Several possible shortcuts remain unproved and are not used. Taking
$y=B$ loses the guaranteed small-prime support. Enforcing divisibility
only by the common growth modulus does not make all beta moduli coprime.
Choosing a generic horizon does not automatically solve the new index CRT
compatibility problem. Dropping $F+t$ from the retained power requests
would remove the stated sign-recovery argument. Treating the new source
as pointwise equivalent to every old witness would also be unjustified.

The next substantive target is to apply the same prime-support principle
to another already authenticated arithmetic quantity, so that another
independent witness can be removed without raising degree. This report
delivers the completed 209-variable construction; no further count is
claimed without a literal polynomial and proof.

\begin{thebibliography}{9}
\bibitem{repo}
Vladimir Reshetnikov, \emph{ProveIt}, inspected snapshot
\nolinkurl{32e9e9eb5999c96b0efe80c80a4018afd33d44ae}.
\url{https://github.com/VladimirReshetnikov/ProveIt/tree/32e9e9eb5999c96b0efe80c80a4018afd33d44ae}.

\bibitem{parent}
\emph{Pell growth from divisibility: an explicit quartic for Con(ZFC)},
incoming package \texttt{ConZFC\_Pell\_Growth\_210\_2026-10-25.zip},
same pinned snapshot. Relevant components are
\texttt{proofs/growth\_proof.md}, \texttt{proofs/AFFINE\_RADIX\_PROOF.md},
\texttt{equations/principal210/equation.json}, and the inherited anchor
package. Exact archive retained in \texttt{inherited/}.
\url{https://github.com/VladimirReshetnikov/ProveIt/blob/32e9e9eb5999c96b0efe80c80a4018afd33d44ae/docs/incoming/ConZFC_Pell_Growth_210_2026-10-25.zip}.

\bibitem{anchor}
\emph{A reused CRT index: an explicit Diophantine consistency equation},
ProveIt source 173, \texttt{173-crt-index-anchor-215}, 25 October 2026.
In particular \texttt{C\_ANCHOR\_THEOREM.md},
\nolinkurl{LABEL_SEPARATION_WITHOUT_DERIVED_RANGE.md},
\texttt{CONSTRUCTOR\_DEPTH.md}, and
\nolinkurl{ORDINARY_PROOF_DEPENDENCY_BOUNDARY.md}.
Selected files and the relevant calibration sources are retained within
the parent archive.
\url{https://github.com/VladimirReshetnikov/ProveIt/tree/32e9e9eb5999c96b0efe80c80a4018afd33d44ae/Computability/DiophantineConsistency/con-zfc-diophantine-forms/173-crt-index-anchor-215}.

\bibitem{linearcrt}
\emph{Linear CRT}, ProveIt source 157, \texttt{157-linear-crt}.
The retained \texttt{MAIN\_PROOF.md} supplies the factorial-free binomial
CRT and signed arithmetic interfaces. See also the retained source-151
factorial-free report and source-158 pooled-power proofs within the parent
archive.
\url{https://github.com/VladimirReshetnikov/ProveIt/tree/32e9e9eb5999c96b0efe80c80a4018afd33d44ae/Computability/DiophantineConsistency/con-zfc-diophantine-forms/157-linear-crt}.

\bibitem{carrier}
\emph{ConZFC Pell Carrier65}, incoming package
\texttt{ConZFC\_Pell\_Carrier65\_2026-10-25.zip}, same pinned snapshot.
Used only to describe the broader comparison frontier, not as a logical
dependency of the new quartic.
\url{https://github.com/VladimirReshetnikov/ProveIt/blob/32e9e9eb5999c96b0efe80c80a4018afd33d44ae/docs/incoming/ConZFC_Pell_Carrier65_2026-10-25.zip}.
\end{thebibliography}

\clearpage
\appendix
\section{The complete literal polynomial}\label{app:residuals}

All variables in the following list range independently over $\Z$.
The polynomial is exactly $P=R_0^2+\cdots+R_{149}^2$. Every residual
has total degree at most two. This list is independently parsed back into
integer coefficient dictionaries and compared with the machine-readable
equation. The source contains no omitted rows, variable exponents, special
function symbols, or externally supplied coefficients.

% Literal quadratic residuals for the 209-variable equation.
\begingroup\small
\setlength{\jot}{2pt}
\begin{align*}
R_{0} ={}& x_{0} - x_{4} + x_{7} - x_{0} x_{4}
\end{align*}
\begin{align*}
R_{1} ={}& -x_{2} - x_{4} x_{5} + x_{6} x_{7}
\end{align*}
\begin{align*}
R_{2} ={}& -1 + x_{1} - x_{6} x_{7}
\end{align*}
\begin{align*}
R_{3} ={}& 1 + x_{2} + 2x_{16} - x_{3} x_{18}
\end{align*}
\begin{align*}
R_{4} ={}& -4 - 3x_{8} - 4x_{9} + x_{19} - x_{8}^{2} \\
&{} -2x_{8} x_{9} - x_{9}^{2}
\end{align*}
\begin{align*}
R_{5} ={}& x_{8} + x_{21} - x_{8}^{2} - 8x_{8} x_{19} - 4x_{8} x_{20} \\
&{} -16x_{19}^{2} - 16x_{19} x_{20} - 4x_{20}^{2}
\end{align*}
\begin{align*}
R_{6} ={}& 1 + 2x_{21} + x_{22} - 4x_{21}^{2}
\end{align*}
\begin{align*}
R_{7} ={}& -2 + x_{23} - x_{9}^{2} - 2x_{9} x_{10} - x_{10}^{2}
\end{align*}
\begin{align*}
R_{8} ={}& 1 + x_{24} - x_{8}^{2} - 2x_{8} x_{10} - x_{10}^{2}
\end{align*}
\begin{align*}
R_{9} ={}& -29 - 61x_{8} + 16x_{9} - 96x_{10} - 28x_{19} \\
&{} + 8x_{23} - 24x_{24} + x_{25} - 35x_{8}^{2} - 96x_{8} x_{10} \\
&{} -24x_{8} x_{19} - 4x_{8} x_{23} - 24x_{8} x_{24} - 2x_{9}^{2} - 8x_{9} x_{23} \\
&{} -64x_{10}^{2} - 32x_{10} x_{19} - 32x_{10} x_{24} - 4x_{19}^{2} - 8x_{19} x_{24} \\
&{} -4x_{23}^{2} - 4x_{24}^{2}
\end{align*}
\begin{align*}
R_{10} ={}& -4 + 10x_{8} + x_{26} - 2x_{43} - 2x_{8}^{2} \\
&{} -8x_{8} x_{9} - 8x_{8} x_{23} - 4x_{25}^{2}
\end{align*}
\begin{align*}
R_{11} ={}& -4 + x_{27} - 16x_{10}^{2} - 16x_{10} x_{23} - 16x_{10} x_{24} \\
&{} -4x_{23}^{2} - 8x_{23} x_{24} - 4x_{24}^{2}
\end{align*}
\begin{align*}
R_{12} ={}& -7 + 24x_{10} + 10x_{19} + 24x_{24} + 12x_{27} \\
&{} + x_{28} - 16x_{10}^{2} - 16x_{10} x_{19} - 32x_{10} x_{24} - 16x_{10} x_{27} \\
&{} -4x_{19}^{2} - 16x_{19} x_{24} - 8x_{19} x_{27} - 16x_{24}^{2} - 16x_{24} x_{27} \\
&{} -4x_{27}^{2}
\end{align*}
\begin{align*}
R_{13} ={}& -114 - 127x_{8} - 67x_{9} - 224x_{10} - 16x_{23} \\
&{} -40x_{24} + x_{29} - 37x_{8}^{2} - 24x_{8} x_{9} - 112x_{8} x_{10} \\
&{} -4x_{8} x_{23} - 24x_{8} x_{24} - 37x_{9}^{2} - 112x_{9} x_{10} - 24x_{9} x_{23} \\
&{} -4x_{9} x_{24} - 128x_{10}^{2} - 32x_{10} x_{23} - 32x_{10} x_{24} - 4x_{23}^{2} \\
&{} -4x_{24}^{2}
\end{align*}
\begin{align*}
R_{14} ={}& -30 - 14x_{8} - 96x_{9} - 128x_{10} - 32x_{23} \\
&{} + x_{30} - 2x_{8}^{2} - 24x_{8} x_{9} - 32x_{8} x_{10} - 8x_{8} x_{23} \\
&{} -72x_{9}^{2} - 192x_{9} x_{10} - 48x_{9} x_{23} - 128x_{10}^{2} - 64x_{10} x_{23} \\
&{} -8x_{23}^{2} - 4x_{29}^{2}
\end{align*}
\begin{align*}
R_{15} ={}& -74 - 42x_{8} - 84x_{9} + 29x_{10} + 21x_{19} \\
&{} + 40x_{23} + x_{31} - x_{59} + 8x_{8} x_{10} + 16x_{8} x_{23} \\
&{} + 16x_{9} x_{10} + 32x_{9} x_{23} - 3x_{10}^{2} - 4x_{10} x_{19} - 8x_{10} x_{23} \\
&{} -8x_{19} x_{23} - 4x_{23}^{2}
\end{align*}
\begin{align*}
R_{16} ={}& -131 - 184x_{8} - 226x_{10} - 88x_{23} - 46x_{24} \\
&{} -44x_{31} + x_{32} - 64x_{8}^{2} - 160x_{8} x_{10} - 64x_{8} x_{23} \\
&{} -32x_{8} x_{24} - 32x_{8} x_{31} - 100x_{10}^{2} - 80x_{10} x_{23} - 40x_{10} x_{24} \\
&{} -40x_{10} x_{31} - 16x_{23}^{2} - 16x_{23} x_{24} - 16x_{23} x_{31} - 4x_{24}^{2} \\
&{} -8x_{24} x_{31} - 4x_{31}^{2}
\end{align*}
\begin{align*}
R_{17} ={}& x_{20} - 4x_{19}^{2}
\end{align*}
\begin{align*}
R_{18} ={}& x_{8} + x_{33} - x_{8}^{2} + 8x_{8} x_{19} - 4x_{8} x_{20} \\
&{} -16x_{19}^{2} + 16x_{19} x_{20} - 4x_{20}^{2}
\end{align*}
\begin{align*}
R_{19} ={}& 1 + 2x_{33} + x_{34} - 4x_{33}^{2}
\end{align*}
\begin{align*}
R_{20} ={}& x_{35} - x_{9}^{2} - 4x_{9} x_{34} - 4x_{34}^{2}
\end{align*}
\begin{align*}
R_{21} ={}& -44 - 52x_{8} - 104x_{9} + 50x_{19} + x_{36} \\
&{} -16x_{8}^{2} - 64x_{8} x_{9} + 32x_{8} x_{19} - 64x_{9}^{2} + 64x_{9} x_{19} \\
&{} -16x_{19}^{2}
\end{align*}
\begin{align*}
R_{22} ={}& 8x_{9} + 6x_{10} + 4x_{36} + x_{37} - 62x_{9}^{2} \\
&{} -92x_{9} x_{10} - 32x_{9} x_{23} - 32x_{9} x_{36} - 34x_{10}^{2} - 24x_{10} x_{23} \\
&{} -24x_{10} x_{36} - 4x_{23}^{2} - 8x_{23} x_{36} - 4x_{36}^{2}
\end{align*}
\begin{align*}
R_{23} ={}& -90 - 152x_{8} - 114x_{10} - 38x_{24} - 36x_{37} \\
&{} + x_{38} - 64x_{8}^{2} - 96x_{8} x_{10} - 32x_{8} x_{24} - 32x_{8} x_{37} \\
&{} -36x_{10}^{2} - 24x_{10} x_{24} - 24x_{10} x_{37} - 4x_{24}^{2} - 8x_{24} x_{37} \\
&{} -4x_{37}^{2}
\end{align*}
\begin{align*}
R_{24} ={}& x_{39} - x_{10}^{2} - 2x_{10} x_{11} - x_{11}^{2}
\end{align*}
\begin{align*}
R_{25} ={}& -7 + 8x_{8} + 6x_{24} + x_{40} - x_{43} \\
&{} -3x_{8}^{2} - 8x_{8} x_{9} - 8x_{8} x_{23} - 8x_{8} x_{24} - 4x_{8} x_{39} \\
&{} -8x_{9} x_{24} - 2x_{10}^{2} - 8x_{10} x_{39} - 8x_{23} x_{24} - 4x_{24}^{2} \\
&{} -4x_{39}^{2}
\end{align*}
\begin{align*}
R_{26} ={}& -4x_{10} - 2x_{24} - 4x_{39} + x_{41} - 16x_{8}^{2} \\
&{} -36x_{8} x_{10} - 32x_{8} x_{24} - 40x_{8} x_{39} - 16x_{8} x_{40} - 22x_{10}^{2} \\
&{} -32x_{10} x_{24} - 48x_{10} x_{39} - 16x_{10} x_{40} - 16x_{24}^{2} - 32x_{24} x_{39} \\
&{} -16x_{24} x_{40} - 24x_{39}^{2} - 16x_{39} x_{40} - 4x_{40}^{2}
\end{align*}
\begin{align*}
R_{27} ={}& x_{10} + x_{42} - 25x_{10}^{2} - 40x_{10} x_{24} - 20x_{10} x_{27} \\
&{} -16x_{24}^{2} - 16x_{24} x_{27} - 4x_{27}^{2}
\end{align*}
\begin{align*}
R_{28} ={}& -1 + 8x_{9} + 8x_{23} + x_{43} - 4x_{9}^{2} \\
&{} -8x_{9} x_{23} - 4x_{23}^{2}
\end{align*}
\begin{align*}
R_{29} ={}& -1 + x_{44} - 16x_{9}^{2} - 8x_{9} x_{10} - 32x_{9} x_{23} \\
&{} -16x_{9} x_{43} - x_{10}^{2} - 8x_{10} x_{23} - 4x_{10} x_{43} - 16x_{23}^{2} \\
&{} -16x_{23} x_{43} - 4x_{43}^{2}
\end{align*}
\begin{align*}
R_{30} ={}& 3 + x_{9} - x_{23} + x_{45} - 2x_{9} x_{10} \\
&{} -4x_{9} x_{44} - 3x_{10}^{2} - 8x_{10} x_{44} - 4x_{44}^{2}
\end{align*}
\begin{align*}
R_{31} ={}& -1 - 4x_{45} + x_{46} - 4x_{45}^{2}
\end{align*}
\begin{align*}
R_{32} ={}& -1 + x_{47} - x_{10}^{2} - 8x_{10} x_{45} - 4x_{10} x_{46} \\
&{} -16x_{45}^{2} - 16x_{45} x_{46} - 4x_{46}^{2}
\end{align*}
\begin{align*}
R_{33} ={}& 2 + x_{9} + x_{48} - x_{9}^{2} - 4x_{9} x_{10} \\
&{} -4x_{9} x_{47} - 4x_{10}^{2} - 8x_{10} x_{47} - 4x_{47}^{2}
\end{align*}
\begin{align*}
R_{34} ={}& 1 + x_{49} - 4x_{42}^{2} - 8x_{42} x_{48} - 4x_{48}^{2}
\end{align*}
\begin{align*}
R_{35} ={}& 1 + x_{8} + x_{50} - x_{8}^{2} - 8x_{8} x_{42} \\
&{} -4x_{8} x_{49} - 16x_{42}^{2} - 16x_{42} x_{49} - 4x_{49}^{2}
\end{align*}
\begin{align*}
R_{36} ={}& -1 - 4x_{50} + x_{51} - 4x_{50}^{2}
\end{align*}
\begin{align*}
R_{37} ={}& -1 + x_{52} - x_{10}^{2} - 8x_{10} x_{50} - 4x_{10} x_{51} \\
&{} -16x_{50}^{2} - 16x_{50} x_{51} - 4x_{51}^{2}
\end{align*}
\begin{align*}
R_{38} ={}& 2 + x_{8} + x_{53} - x_{8}^{2} - 4x_{8} x_{10} \\
&{} -4x_{8} x_{52} - 4x_{10}^{2} - 8x_{10} x_{52} - 4x_{52}^{2}
\end{align*}
\begin{align*}
R_{39} ={}& -10 - 12x_{53} + x_{54} - 4x_{53}^{2}
\end{align*}
\begin{align*}
R_{40} ={}& -19 + 4x_{9} + 4x_{23} - 16x_{53} + 6x_{54} \\
&{} + x_{55} - 4x_{9}^{2} - 8x_{9} x_{23} - 16x_{9} x_{53} - 8x_{9} x_{54} \\
&{} -4x_{23}^{2} - 16x_{23} x_{53} - 8x_{23} x_{54} - 24x_{53}^{2} - 16x_{53} x_{54} \\
&{} -4x_{54}^{2}
\end{align*}
\begin{align*}
R_{41} ={}& -10 + x_{9} - 12x_{48} + x_{56} - 25x_{9}^{2} \\
&{} -40x_{9} x_{23} - 20x_{9} x_{55} - 16x_{23}^{2} - 16x_{23} x_{55} - 4x_{48}^{2} \\
&{} -4x_{55}^{2}
\end{align*}
\begin{align*}
R_{42} ={}& 21 + 24x_{48} - 2x_{56} + x_{57} - 8x_{48}^{2} \\
&{} -16x_{48} x_{56} - 4x_{56}^{2}
\end{align*}
\begin{align*}
R_{43} ={}& -1 + 4x_{8} + 8x_{9} + x_{10} - 2x_{19} \\
&{} + x_{58} + 12x_{59} - x_{10}^{2} - 4x_{10} x_{57} - 4x_{57}^{2} \\
&{} -4x_{59}^{2}
\end{align*}
\begin{align*}
R_{44} ={}& -22 - 38x_{8} - 76x_{9} + 11x_{10} + 19x_{19} \\
&{} + x_{59} - 16x_{8}^{2} - 64x_{8} x_{9} + 8x_{8} x_{10} + 16x_{8} x_{19} \\
&{} -64x_{9}^{2} + 16x_{9} x_{10} + 32x_{9} x_{19} - x_{10}^{2} - 4x_{10} x_{19} \\
&{} -4x_{19}^{2}
\end{align*}
\begin{align*}
R_{45} ={}& -4 + 8x_{8} + 16x_{9} - 4x_{19} + 24x_{59} \\
&{} + x_{60} - 4x_{58}^{2} - 8x_{59}^{2}
\end{align*}
\begin{align*}
R_{46} ={}& -1 - 4x_{60} + x_{61} - 4x_{60}^{2}
\end{align*}
\begin{align*}
R_{47} ={}& x_{62} - x_{2} x_{13}
\end{align*}
\begin{align*}
R_{48} ={}& x_{12} - x_{13} - x_{63} - 2x_{22} x_{62}
\end{align*}
\begin{align*}
R_{49} ={}& x_{12} - x_{13} - x_{64} - 2x_{26} x_{62}
\end{align*}
\begin{align*}
R_{50} ={}& x_{12} - x_{13} - 2x_{62} - x_{65} - 2x_{28} x_{62}
\end{align*}
\begin{align*}
R_{51} ={}& x_{12} - x_{13} - x_{66} - 2x_{30} x_{62}
\end{align*}
\begin{align*}
R_{52} ={}& x_{12} - x_{13} + 2x_{62} - x_{67} - 2x_{32} x_{62}
\end{align*}
\begin{align*}
R_{53} ={}& x_{12} - x_{13} - x_{68} - 4x_{34} x_{62} - 2x_{35} x_{62}
\end{align*}
\begin{align*}
R_{54} ={}& x_{12} - x_{13} - x_{69} - 2x_{38} x_{62}
\end{align*}
\begin{align*}
R_{55} ={}& -2x_{62} + x_{65} - x_{70} + 2x_{28} x_{62} - 2x_{41} x_{62}
\end{align*}
\begin{align*}
R_{56} ={}& x_{71} - x_{8} x_{9} + x_{8} x_{10} + x_{9} x_{10} - x_{10}^{2}
\end{align*}
\begin{align*}
R_{57} ={}& 1 + 2x_{11} + x_{72} + x_{11}^{2} - x_{71}^{2}
\end{align*}
\begin{align*}
R_{58} ={}& x_{73} - x_{12}^{2} + 2x_{12} x_{13} + 2x_{12} x_{74} - x_{13}^{2} \\
&{} -2x_{13} x_{74} - x_{72}^{2} - x_{74}^{2}
\end{align*}
\begin{align*}
R_{59} ={}& x_{74} - 4x_{60} x_{62} - 2x_{61} x_{62}
\end{align*}
\begin{align*}
R_{60} ={}& -x_{13} + x_{77} - x_{9} x_{62} + x_{75} x_{76}
\end{align*}
\begin{align*}
R_{61} ={}& -1 + 8x_{2} + x_{75} + 4x_{2} x_{8} + 8x_{2} x_{9} \\
&{} -2x_{2} x_{19}
\end{align*}
\begin{align*}
R_{62} ={}& -x_{12} + x_{76} - x_{8} x_{201}
\end{align*}
\begin{align*}
R_{63} ={}& x_{67} + x_{78} + x_{9} x_{201} - 2x_{19} x_{62} + 2x_{32} x_{62}
\end{align*}
\begin{align*}
R_{64} ={}& x_{79} - x_{12} x_{63} + x_{13} x_{63}
\end{align*}
\begin{align*}
R_{65} ={}& x_{80} - x_{64} x_{79}
\end{align*}
\begin{align*}
R_{66} ={}& x_{81} - x_{65} x_{80}
\end{align*}
\begin{align*}
R_{67} ={}& x_{82} - x_{66} x_{81}
\end{align*}
\begin{align*}
R_{68} ={}& x_{83} - x_{67} x_{82}
\end{align*}
\begin{align*}
R_{69} ={}& x_{84} - x_{68} x_{83}
\end{align*}
\begin{align*}
R_{70} ={}& x_{85} - x_{69} x_{84}
\end{align*}
\begin{align*}
R_{71} ={}& x_{86} - x_{70} x_{85}
\end{align*}
\begin{align*}
R_{72} ={}& x_{87} - x_{77} x_{86}
\end{align*}
\begin{align*}
R_{73} ={}& x_{88} - x_{78} x_{87}
\end{align*}
\begin{align*}
R_{74} ={}& x_{89} - x_{88}^{2}
\end{align*}
\begin{align*}
R_{75} ={}& -x_{4} + 2x_{16} + x_{90} - 2x_{4} x_{16}
\end{align*}
\begin{align*}
R_{76} ={}& x_{91} - x_{14} x_{90}
\end{align*}
\begin{align*}
R_{77} ={}& -2x_{1}^{2} + 2x_{1} x_{12} + 2x_{1} x_{13} + 2x_{1} x_{91} + 2x_{1} x_{92} \\
&{} -x_{12}^{2} - 2x_{12} x_{91} - x_{13}^{2} - 2x_{13} x_{92} + x_{17} x_{93} \\
&{} -x_{73} x_{89} - x_{91}^{2} - x_{92}^{2}
\end{align*}
\begin{align*}
R_{78} ={}& x_{15} + x_{92} - x_{4} x_{15} - x_{15} x_{90}
\end{align*}
\begin{align*}
R_{79} ={}& x_{97} + x_{99} - x_{94} x_{97}
\end{align*}
\begin{align*}
R_{80} ={}& -x_{98} - x_{3} x_{120} + x_{95} x_{193} - x_{99} x_{120}
\end{align*}
\begin{align*}
R_{81} ={}& -x_{102} + x_{100} x_{193} - x_{101} x_{123}
\end{align*}
\begin{align*}
R_{82} ={}& -x_{105} + x_{103} x_{193} - x_{104} x_{123}
\end{align*}
\begin{align*}
R_{83} ={}& -x_{108} + x_{106} x_{193} - x_{107} x_{123}
\end{align*}
\begin{align*}
R_{84} ={}& -x_{111} + x_{109} x_{193} - x_{110} x_{123}
\end{align*}
\begin{align*}
R_{85} ={}& -x_{115} + x_{112} x_{194} - x_{114} x_{123}
\end{align*}
\begin{align*}
R_{86} ={}& -x_{118} + x_{116} x_{194} - x_{117} x_{123}
\end{align*}
\begin{align*}
R_{87} ={}& x_{123} - x_{17} x_{120}
\end{align*}
\begin{align*}
R_{88} ={}& x_{119} - x_{122} - x_{121} x_{123}
\end{align*}
\begin{align*}
R_{89} ={}& x_{124} - x_{126} - x_{123} x_{125}
\end{align*}
\begin{align*}
R_{90} ={}& -x_{129} - x_{123} x_{128} + x_{127} x_{193}
\end{align*}
\begin{align*}
R_{91} ={}& x_{130} + x_{133} - 2x_{130} x_{132}
\end{align*}
\begin{align*}
R_{92} ={}& -x_{133} + x_{134} - x_{130} x_{133}
\end{align*}
\begin{align*}
R_{93} ={}& -1 + x_{131}^{2} + x_{133}^{2} - x_{134}^{2}
\end{align*}
\begin{align*}
R_{94} ={}& 1 - x_{135} + x_{137} - x_{131} x_{136}
\end{align*}
\begin{align*}
R_{95} ={}& -1 + 2x_{135} - x_{135}^{2} + 2x_{135} x_{137} + x_{136}^{2}
\end{align*}
\begin{align*}
R_{96} ={}& 1 + x_{138} + x_{3}^{2} - 2x_{3} x_{131}
\end{align*}
\begin{align*}
R_{97} ={}& x_{17} - x_{169} - x_{3} x_{170} + x_{138} x_{139}
\end{align*}
\begin{align*}
R_{98} ={}& -5 - 8x_{2} - 16x_{16} + 4x_{131} - x_{140} \\
&{} -4x_{2}^{2} - 16x_{2} x_{16} + 4x_{2} x_{131} - 16x_{16}^{2} + 8x_{16} x_{131}
\end{align*}
\begin{align*}
R_{99} ={}& -1 + x_{94} - x_{169} - 2x_{170} - 2x_{2} x_{170} \\
&{} -4x_{16} x_{170} + x_{140} x_{141}
\end{align*}
\begin{align*}
R_{100} ={}& -x_{142} + x_{144} - x_{131} x_{143}
\end{align*}
\begin{align*}
R_{101} ={}& -1 - x_{142}^{2} + 2x_{142} x_{144} + x_{143}^{2}
\end{align*}
\begin{align*}
R_{102} ={}& -x_{145} + x_{147} - x_{131} x_{146}
\end{align*}
\begin{align*}
R_{103} ={}& -1 - x_{145}^{2} + 2x_{145} x_{147} + x_{146}^{2}
\end{align*}
\begin{align*}
R_{104} ={}& -x_{148} + x_{150} - x_{131} x_{149}
\end{align*}
\begin{align*}
R_{105} ={}& -1 - x_{148}^{2} + 2x_{148} x_{150} + x_{149}^{2}
\end{align*}
\begin{align*}
R_{106} ={}& -x_{151} + x_{153} - x_{131} x_{152}
\end{align*}
\begin{align*}
R_{107} ={}& -1 - x_{151}^{2} + 2x_{151} x_{153} + x_{152}^{2}
\end{align*}
\begin{align*}
R_{108} ={}& -x_{154} + x_{156} - x_{131} x_{155}
\end{align*}
\begin{align*}
R_{109} ={}& -1 - x_{154}^{2} + 2x_{154} x_{156} + x_{155}^{2}
\end{align*}
\begin{align*}
R_{110} ={}& -x_{157} + x_{159} - x_{131} x_{158}
\end{align*}
\begin{align*}
R_{111} ={}& -1 - x_{157}^{2} + 2x_{157} x_{159} + x_{158}^{2}
\end{align*}
\begin{align*}
R_{112} ={}& 1 - x_{160} + x_{162} - x_{131} x_{161}
\end{align*}
\begin{align*}
R_{113} ={}& -1 + 2x_{160} - x_{160}^{2} + 2x_{160} x_{162} + x_{161}^{2}
\end{align*}
\begin{align*}
R_{114} ={}& -x_{163} + x_{165} - x_{131} x_{164}
\end{align*}
\begin{align*}
R_{115} ={}& -1 - x_{163}^{2} + 2x_{163} x_{165} + x_{164}^{2}
\end{align*}
\begin{align*}
R_{116} ={}& -x_{166} + x_{168} - x_{131} x_{167}
\end{align*}
\begin{align*}
R_{117} ={}& -1 - x_{166}^{2} + 2x_{166} x_{168} + x_{167}^{2}
\end{align*}
\begin{align*}
R_{118} ={}& 1 - x_{169} + x_{171} - x_{131} x_{170}
\end{align*}
\begin{align*}
R_{119} ={}& -1 + 2x_{169} - x_{169}^{2} + 2x_{169} x_{171} + x_{170}^{2}
\end{align*}
\begin{align*}
R_{120} ={}& -x_{172} + x_{174} - x_{131} x_{173}
\end{align*}
\begin{align*}
R_{121} ={}& -1 - x_{172}^{2} + 2x_{172} x_{174} + x_{173}^{2}
\end{align*}
\begin{align*}
R_{122} ={}& -x_{175} + x_{177} - x_{131} x_{176}
\end{align*}
\begin{align*}
R_{123} ={}& -1 - x_{175}^{2} + 2x_{175} x_{177} + x_{176}^{2}
\end{align*}
\begin{align*}
R_{124} ={}& -8 - 2x_{0} - 8x_{2} - 16x_{16} + 2x_{94} \\
&{} + x_{130} - 4x_{0}^{2} - 2x_{0} x_{2} - x_{1}^{2} - 6x_{2}^{2} \\
&{} -16x_{2} x_{16} - x_{3}^{2} - x_{4}^{2} - x_{8}^{2} - x_{9}^{2} \\
&{} -x_{10}^{2} - x_{11}^{2} - x_{12}^{2} - x_{13}^{2} - x_{14}^{2} \\
&{} -x_{15}^{2} - 20x_{16}^{2} - x_{17}^{2} - 2x_{94}^{2} - x_{95}^{2} \\
&{} -x_{96}^{2} - x_{100}^{2} - x_{103}^{2} - x_{106}^{2} - x_{109}^{2} \\
&{} -x_{112}^{2} - x_{113}^{2} - x_{116}^{2} - x_{119}^{2} - x_{120}^{2} \\
&{} -x_{124}^{2} - x_{127}^{2}
\end{align*}
\begin{align*}
R_{125} ={}& -32 + x_{0} - 128x_{178} - 128x_{178}^{2}
\end{align*}
\begin{align*}
R_{126} ={}& x_{189} - x_{131} x_{180}
\end{align*}
\begin{align*}
R_{127} ={}& -1 + x_{179}^{2} + x_{180}^{2} - x_{189}^{2}
\end{align*}
\begin{align*}
R_{128} ={}& x_{190} - x_{182} x_{183}
\end{align*}
\begin{align*}
R_{129} ={}& -1 + x_{181}^{2} + x_{182}^{2} - x_{190}^{2}
\end{align*}
\begin{align*}
R_{130} ={}& x_{191} - x_{176}^{2}
\end{align*}
\begin{align*}
R_{131} ={}& x_{180} - x_{191} - 2x_{185} x_{191}
\end{align*}
\begin{align*}
R_{132} ={}& -1 - 4x_{180} + x_{183} - 8x_{180} x_{184}
\end{align*}
\begin{align*}
R_{133} ={}& -x_{131} + x_{183} - x_{179} x_{186}
\end{align*}
\begin{align*}
R_{134} ={}& -x_{177} + x_{181} - x_{179} x_{187}
\end{align*}
\begin{align*}
R_{135} ={}& -2x_{16} + x_{182} - 4x_{176} x_{188}
\end{align*}
\begin{align*}
R_{136} ={}& 9 - 2x_{144} - 2x_{147} - 2x_{150} - 2x_{153} \\
&{} -2x_{156} - 2x_{159} - 2x_{165} - 2x_{168} - 2x_{174} \\
&{} -2x_{177} + x_{192} - x_{197} + 2x_{0}^{2} + 2x_{0} x_{2} \\
&{} -2x_{0} x_{136} - 2x_{0} x_{143} + x_{1}^{2} - 2x_{1} x_{170} + x_{2}^{2} \\
&{} -2x_{2} x_{143} + x_{4}^{2} - 2x_{4} x_{161} + x_{8}^{2} - 2x_{8} x_{146} \\
&{} + x_{9}^{2} - 2x_{9} x_{149} + x_{10}^{2} - 2x_{10} x_{152} + x_{11}^{2} \\
&{} -2x_{11} x_{155} + x_{12}^{2} - 2x_{12} x_{158} + x_{13}^{2} - 2x_{13} x_{164} \\
&{} + x_{14}^{2} - 2x_{14} x_{167} + x_{15}^{2} - 2x_{15} x_{173} - 4x_{16} x_{176} \\
&{} -x_{131} x_{192} + x_{137}^{2} + x_{144}^{2} + x_{147}^{2} + x_{150}^{2} \\
&{} + x_{153}^{2} + x_{156}^{2} + x_{159}^{2} + x_{162}^{2} + x_{165}^{2} \\
&{} + x_{168}^{2} + x_{171}^{2} + x_{174}^{2} + 2x_{176}^{2} + x_{177}^{2} \\
&{} -x_{195}^{2} - x_{196}^{2} - x_{197}^{2}
\end{align*}
\begin{align*}
R_{137} ={}& x_{120} - x_{96} x_{193}
\end{align*}
\begin{align*}
R_{138} ={}& x_{120} - x_{113} x_{194}
\end{align*}
\begin{align*}
R_{139} ={}& -1 - x_{197} - 4x_{16}^{2} - x_{136}^{2} - x_{143}^{2} \\
&{} -x_{146}^{2} - x_{149}^{2} - x_{152}^{2} - x_{155}^{2} - x_{158}^{2} \\
&{} -x_{161}^{2} - x_{164}^{2} - x_{167}^{2} - x_{170}^{2} - x_{173}^{2} \\
&{} + x_{176}^{2} - x_{195}^{2} - x_{196}^{2} - x_{197}^{2}
\end{align*}
\begin{align*}
R_{140} ={}& -3 + 3x_{200} - 4x_{98}^{2} - 4x_{102}^{2} - 4x_{105}^{2} \\
&{} -4x_{108}^{2} - 4x_{111}^{2} - 4x_{115}^{2} - 4x_{118}^{2} + x_{120}^{2} \\
&{} -4x_{122}^{2} - 4x_{126}^{2} - 4x_{129}^{2} - x_{198}^{2} - x_{199}^{2} \\
&{} -x_{200}^{2}
\end{align*}
\begin{align*}
R_{141} ={}& x_{201} - x_{2} x_{12}
\end{align*}
\begin{align*}
R_{142} ={}& x_{202} + x_{94} x_{203} - 2x_{95} x_{143} - 2x_{100} x_{146} - 2x_{103} x_{149} \\
&{} -2x_{106} x_{152} - 2x_{109} x_{155} - 2x_{112} x_{158} - 2x_{116} x_{164} - 2x_{119} x_{167} \\
&{} -2x_{124} x_{173} - 2x_{127} x_{176} - 2x_{131} x_{203} + 2x_{143} x_{144} + 2x_{146} x_{147} \\
&{} + 2x_{149} x_{150} + 2x_{152} x_{153} + 2x_{155} x_{156} + 2x_{158} x_{159} + 2x_{164} x_{165} \\
&{} + 2x_{167} x_{168} + 2x_{173} x_{174} + 2x_{176} x_{177}
\end{align*}
\begin{align*}
R_{143} ={}& -10 + x_{203} - x_{94} x_{202} + x_{95}^{2} - 2x_{95} x_{142} \\
&{} + x_{100}^{2} - 2x_{100} x_{145} + x_{103}^{2} - 2x_{103} x_{148} + x_{106}^{2} \\
&{} -2x_{106} x_{151} + x_{109}^{2} - 2x_{109} x_{154} + x_{112}^{2} - 2x_{112} x_{157} \\
&{} + x_{116}^{2} - 2x_{116} x_{163} + x_{119}^{2} - 2x_{119} x_{166} + x_{124}^{2} \\
&{} -2x_{124} x_{172} + x_{127}^{2} - 2x_{127} x_{175} + 2x_{142} x_{144} + 2x_{145} x_{147} \\
&{} + 2x_{148} x_{150} + 2x_{151} x_{153} + 2x_{154} x_{156} + 2x_{157} x_{159} + 2x_{163} x_{165} \\
&{} + 2x_{166} x_{168} + 2x_{172} x_{174} + 2x_{175} x_{177}
\end{align*}
\begin{align*}
R_{144} ={}& 2x_{136} + 2x_{161} + 2x_{170} + x_{204} - x_{205} \\
&{} + x_{94} x_{205} - 2x_{96} x_{136} - 2x_{113} x_{161} - 2x_{120} x_{170} - 2x_{131} x_{205} \\
&{} + 2x_{136} x_{137} + 2x_{161} x_{162} + 2x_{170} x_{171}
\end{align*}
\begin{align*}
R_{145} ={}& -2 + 2x_{135} + 2x_{160} + x_{204} + x_{205} \\
&{} -x_{94} x_{204} + x_{96}^{2} - 2x_{96} x_{135} + x_{113}^{2} - 2x_{113} x_{160} \\
&{} + x_{120}^{2} - 2x_{120} x_{169} + 2x_{135} x_{137} + 2x_{160} x_{162} + x_{169}^{2} \\
&{} -x_{170}^{2}
\end{align*}
\begin{align*}
R_{146} ={}& -1 + x_{4} - x_{2} x_{3}
\end{align*}
\begin{align*}
R_{147} ={}& x_{207} - 524309x_{0} x_{2}
\end{align*}
\begin{align*}
R_{148} ={}& -1 - 2x_{2} x_{207} + x_{206}^{2} - x_{207}^{2}
\end{align*}
\begin{align*}
R_{149} ={}& x_{0} + x_{2} - 524309x_{0} x_{208}
\end{align*}
\endgroup


\clearpage
\section{Artifact map and integration notes}

\begin{longtable}{p{0.38\textwidth}p{0.56\textwidth}}
\toprule
Path & Purpose\\
\midrule
\endhead
\texttt{equations/equation.json} & Complete quadratic residual system defining the one quartic.\\
\nolinkurl{equations/expanded_polynomial.json} & All collected integer coefficients of $P$.\\
\texttt{equations/residuals.txt} & Readable literal residual list.\\
\texttt{equations/variables.tsv} & All 209 coordinate names and source roles.\\
\texttt{equations/transformation.json} & Exact changed rows, parent hash, and index map.\\
\texttt{code/build.py} & Deterministic standard-library producer.\\
\texttt{code/verify.py} & Independent exact coefficient verifier.\\
\texttt{code/check\_printed.py} & Parser verifying the printed polynomial.\\
\texttt{code/check\_arithmetic.py} & Independent modular arithmetic regression.\\
\texttt{code/evaluate.py} & Exact evaluation on a supplied integer tuple.\\
\texttt{proofs/} & Independent mathematical review notes.\\
\texttt{verification/} & Fresh finite verification receipts, including the parent replay.\\
\texttt{inherited/} & Exact P210 parent archive, its equation bytes, and license.\\
\texttt{PROVENANCE.json} & Snapshot, source identities, and file hashes.\\
\texttt{SHA256SUMS} & Delivered-file integrity manifest.\\
\bottomrule
\end{longtable}

Suggested additive placement is a new source directory under
\nolinkurl{Computability/DiophantineConsistency/con-zfc-diophantine-forms/},
with a new source number assigned by the repository maintainer. The archive
is suitable for the \texttt{docs/incoming/} intake flow. Existing reports
and their historical metrics should be preserved. This release has not
written to the remote repository.

\end{document}
